\begin{abstract}
Every time series forecasting model consumes a fixed-length history window, yet the window length is chosen by convention rather than by evidence. We show that this convention is silently expensive. Under a protocol that holds the training data budget fixed, the optimal lookback varies by 30x across benchmarks: exchange-rate forecasting is best at 96 steps and degrades 3.6x with longer inputs, while electricity and traffic keep improving up to 2880 steps. The pattern replicates without any training in a zero-shot foundation model, so it is a property of the data--model pairing, not of a training procedure. We explain the divergence with a two-axis account: distant history benefits through periodic and long-memory structure that survives regime changes, and harms through non-stationary drift that turns old windows into contaminated samples. The natural single-axis alternative---the span equals the expected run length of a changepoint model---is falsified: the most regime-switching datasets have the longest optimal lookbacks. The effective span is partially predictable from cheap, training-free statistics of the training data alone (leave-one-dataset-out Spearman $|\rho| = 0.63$ over fourteen public datasets), and the failure of standard architectures to suppress stale content is mechanical: attention flattens with context length, and even a linear model does not learn to gate out distant lags. On controlled synthetic series the harm axis obeys an empirical law, $W^* \propto \lambda^{-1/3}$ in the drift rate $\lambda$, which refutes the classical $-2/3$ exponent of continuous-drift estimation and prescribes a drift-matched exponential input decay; on real benchmarks the calibrated decay saturates to a universal safe default that never degrades the no-budget baseline. A hard input budget recovers 74--98\% of the long-context degradation on trained patch-family models, and budgeted truncation of a frozen foundation model matches validation-based span selection on all eight benchmarks tested. Our results argue that lookback should be treated as a measurable, predictable, data-dependent resource rather than a convention.
\end{abstract}
